\section{Algebraic entropy: uniform gaps and finite-set attainment}
\label{sec:entropy}

The algebraic-entropy strategy replaces the scalar question by a uniform gap for discrete abelian flows and by finite-set attainment.
Its subsidiary directions seek additional control of actual power orbits through Fourier estimates, or a finite-capture mechanism through neutral completion and graph compression.
Both directions must retain unrestricted group and rank.

\begin{definition}\label{fo:algebraic-entropy}
For an endomorphism $\phi$ of a discrete abelian group $G$ and a nonempty finite subset $F\subset G$, put
\[
 T_n(\phi,F)=F+\phi(F)+\cdots+\phi^{n-1}(F),\qquad
 H(\phi,F)=\lim_{n\to\infty}\frac{\log|T_n(\phi,F)|}{n},
\]
and define $h_{\rm alg}(\phi)=\sup_F H(\phi,F)$.
\end{definition}
Subadditivity gives the limit.
This is the finite-subset algebraic entropy of \citet{DGB16}; the sets $F$ need not be subgroups.

\begin{proposition}\label{fo:entropy-bridge}
The infimum of the positive values of this entropy over all abelian groups and endomorphisms equals~\eqref{fo:gap} \cite[Theorem~1.5]{DGB16}.
\end{proposition}

\begin{proposition}\label{fo:yuzvinski-formula}
For a rational linear map $A$ on $\Q^d$, the rational algebraic Yuzvinski formula is
\begin{equation}\label{fo:yuzvinski}
 h_{\rm alg}(A)=m(\chi_{A,\Z}),
\end{equation}
where $\chi_{A,\Z}$ is the primitive integer associate, with positive leading coefficient, of the rational characteristic polynomial \cite[Theorem~4.17 in the consulted preprint]{GBV15}.
\end{proposition}
Clearing denominators contributes to Mahler measure; a sum over expanding archimedean eigenvalues alone would omit that contribution.

The associated attainment question concerns a finite set in the same flow, rather than another flow with the same entropy value.
\begin{proposition}\label{fo:attainment-implication}
If $\mathfrak L>0$, every discrete abelian flow $(G,\phi)$ of finite algebraic entropy has a nonempty finite $F\subset G$ with $H(\phi,F)=h_{\rm alg}(\phi)$ \cite[Theorem~1.7]{DGB16}.
\end{proposition}
A direct construction explains the opposite obstruction without needing a description of the entire entropy spectrum: if positive Mahler measures tend to zero, choose distinct irreducible companion flows with summable positive entropies.
Their direct sum has finite positive entropy, but every finite set lies in finitely many summands and misses a positive tail.
The addition and directed-continuity properties used here are those of \citet{DGB16}.

\begin{proposition}\label{fo:primary}
Let a rational vector space $V$, with endomorphism $A$, have finite algebraic entropy.
Regard $V$ as a $\Q[t]$-module by $t\cdot v=Av$.
For every monic irreducible $p\in\Q[t]$, let
\[
 V_p=\{v:p(A)^jv=0\text{ for some }j\ge1\},\quad
 d_p=\deg p,\quad w_p=m(p_{\Z}).
\]
Here $p_{\Z}$ is the primitive integer associate of $p$ with positive leading coefficient.
Then $V=\bigoplus_pV_p$ and
\begin{equation}\label{fo:primary-sum}
 h_{\rm alg}(A)=\sum_{p:w_p>0}\ell_pw_p,
 \qquad \ell_p=\dim_{\Q}(V_p)/d_p.
\end{equation}
Each positive-weight multiplicity is a finite integer.
Finite-set attainment holds exactly when only finitely many positive-weight primary components occur.
Zero-weight components may have infinite dimension.
\end{proposition}
\begin{proof}
A vector with linearly independent forward iterates supports an integer shift, whose finite alphabets have arbitrarily large entropy.
Monotonicity therefore makes every vector polynomially torsion.
Bezout identities split each vector into the primary components of an annihilating polynomial.
A finite set consequently generates a finite-dimensional invariant rational space.
Formula~\eqref{fo:yuzvinski}, additivity and directed continuity yield~\eqref{fo:primary-sum}.
An infinite positive-weight multiplicity would give infinite entropy.
If only finitely many positive components occur, their finite-dimensional sum captures all the entropy and the finite-dimensional attainment theorem applies.
Otherwise every finite set misses at least one positive component.
\end{proof}

The proposition rules out obtaining nonattainment merely by increasing Jordan lengths at one positive-weight polynomial.
It does not exclude a summable family of distinct positive Mahler weights.
That arithmetic question retains the difficulty of the original problem.


\subsection{Finite-window Fourier budgets and actual transition amplitudes}
\label{sec:fourier}

This subsidiary strategy compares a complete finite-window Fourier contribution with the small entropy deficit already proved for the actual power-orbit law.
The probability law must be specified before one can use that deficit; it is different from algebraic entropy.

\begin{definition}\label{fo:binary-law}
Let $\alpha\ne0$ be an algebraic integer, write $P$ for its minimal polynomial, $d=\deg P$ and $m=m(P)$, and let $\varepsilon_j$ be independent, each taking the values $0,1$ with probability $1/2$.
Define
\[
 S_n=\sum_{j=0}^{n-1}\varepsilon_j\alpha^j,
 \qquad \mu_n(x)=\Prob(S_n=x).
\]
For its finite distribution use the support, Shannon and collision entropies, in the standard R\'enyi conventions \cite[p.~1]{Li26},
\[
 H_0(n)=\log|\supp\mu_n|,\quad
 H_1(n)=-\sum_x\mu_n(x)\log\mu_n(x),\quad
 H_2(n)=-\log\sum_x\mu_n(x)^2.
\]
\end{definition}
The limits $h_j=\lim_n H_j(n)/n$ exist by block subadditivity.
These are entropies of one specified probability law, distinct from $h_{\rm alg}$.
Let $\eta_j=\varepsilon_j-\varepsilon'_j$, where the primed family is an independent copy.
The collision probability is an actual return probability:
\begin{equation}\label{fo:collision-return}
 p_n:=\sum_x\mu_n(x)^2
 =\Prob\left(\sum_{j=0}^{n-1}\eta_j\alpha^j=0\right),
 \qquad \Prob(\eta_j=-1,0,1)=(1/4,1/2,1/4).
\end{equation}
Thus $h_2=-\lim_n n^{-1}\log p_n$.
Changing this original digit law or counting only selected return words changes the object whose rate must be controlled.

\subsubsection{A quantitative baseline retained from the earlier work}
\begin{theorem}\label{fo:convolution-comparison}
For $k\ge1$, replace each digit by the sum of $k$ independent fair signs, and denote the Shannon rate by $h_1^{[k]}$.
The following comparison holds for every nonzero algebraic integer, including those with unit-circle conjugates.
Put
\[
 a=e^{2m}-1,\qquad C(m)=\frac{5a^2}{12}\bigl(1-(1+a)^{-2}\bigr).
\]
Then
\begin{equation}\label{fo:convolution-rate}
 0\le m-h_1^{[k]}\le\frac{C(m)}{k^2},\qquad
 h_1^{[k]}\nearrow m,\qquad
 C(m)=\frac{20}{3}m^3+O(m^4)\quad(m\to0).
\end{equation}
\end{theorem}
For $k=1$ this is the original binary Shannon rate.
For clarity, the convolved digit distribution is not uniform on its alphabet.

\begin{proof}[Proof outline]
The mechanism behind the quantitative statement is as follows.
Represent $1,\alpha,\ldots,\alpha^{n-1}$ in a rational power basis, with coefficient columns $v_j$, and let $W_{n-1}=\sum_{j<n-1}v_jv_j^T$.
For $n>d$ the endpoint quantity $a_n=v_{n-1}^TW_{n-1}^{-1}v_{n-1}$ satisfies $a_n\ge e^{2m}-1$ by a finite Gram determinant and polynomial Jensen argument.
Revealing the endpoint signs successively leaves the exact leverage values $a_n/(k+ra_n)$, $1\le r\le k$.
For a fair-sign linear observation, the posterior second- and fourth-moment bounds give mutual information at least
\[
 F(\ell)=\frac\ell2+\frac{\ell^2(3-2\ell)^4}{324}.
\]
The chain rule therefore bounds each limiting rate below by $\sum_{r=1}^kF(a/(k+ra))$.
Summing the resulting rational remainder estimate gives~\eqref{fo:convolution-rate}.
Independent convolution increases entropy, while the norm-grid bound supplies the upper bound $m$.
No limiting posterior distribution or hyperbolicity assumption is needed.
\end{proof}

\begin{remark}
Closeness to the Mahler scale does not separate that scale from zero.
Equation~\eqref{fo:convolution-rate} is compatible with both $m\to0$ and $h_1\to0$.
A new entropy argument must produce an additional uniform constraint on the actual power-orbit law.
Strictness at a finite time, or a generic posterior identity detached from that orbit, is insufficient.
\end{remark}
\subsubsection{A fixed-input transition obstruction}
\label{fo:transition}

The following result limits a particular single-multiplier Fourier construction.
Fix a nontorsion algebraic unit $\alpha$ with minimal polynomial $P$ and at least one unit-circle conjugate.
Put $K=\Q(\alpha)$ and let $\mathcal O_K$ be its ring of algebraic integers.
For $\theta\in\mathcal O_K\setminus\{0\}$ and $j\in\Z$, define the stable, expanding and central weighted traces
\[
 u_j=\sum_{|\sigma(\alpha)|<1}\sigma(\theta)\sigma(\alpha)^j,\quad
 v_j=\sum_{|\sigma(\alpha)|>1}\sigma(\theta)\sigma(\alpha)^j,\quad
 w_j=\sum_{|\sigma(\alpha)|=1}\sigma(\theta)\sigma(\alpha)^j.
\]
Every embedding is counted, including both members of a complex pair, so these traces are real and $u_j+v_j+w_j$ is an integer.
Call $J\ge0$ admissible if $|u_{J+r}|,|v_{-J-r}|\le1/16$ for every integer $r\ge1$.

\begin{proposition}\label{fo:transition-collapse}
For every sequence $\theta_n\in\mathcal O_K\setminus\{0\}$ whose central conjugates all tend to zero, and every sequence of admissible integers $J_n$, one has
\[
 \prod_{j=-J_n}^{-1}|\cos(2\pi v_j)|
 \prod_{j=0}^{J_n}|\cos(2\pi(u_j+w_j))|
 \longrightarrow0,
\]
where the traces in the $n$th product use $\theta_n$.
No height bound on $\theta_n$ or prescribed growth of $J_n$ is required.
\end{proposition}
\begin{proof}[Proof outline]
Suppose the products have a positive lower bound.
Admissibility and the fixed Vandermonde matrices of the hyperbolic conjugates give uniformly square-summable tails.
On the core, the inequality $-\log\cos(\pi x)\ge\pi^2x^2/2$ bounds the squared nearest-integer errors.
Removing the vanishing central tail leaves a bounded $\ell^2(\Z)$ sequence.
Its nonzero integer defect under the recurrence $P$ can be translated to the origin.
A diagonal limit then gives a decaying response strictly inside the nearest-integer cube with a nonzero integer defect, contradicting the fixed-polynomial integer-defect rigidity used in the complete argument.
The translation includes defects near the window boundary and does not require spatial tightness.
\end{proof}
This is a fixed-$P$ statement without an effective decay modulus.
It leaves changed finite windows, jointly chosen multipliers and varying-input estimates as separate questions.

\subsubsection{A conditional critical regularity theorem}
\label{ret:fractional}

Return to an algebraic unit $\alpha$ and put $\mathcal O=\Z[\alpha]$.
For the original walk $X_0=0$, $X_{n+1}=\alpha X_n+\eta_n$, let $f_n$ be the complete first positive return probabilities and put
\[
 R=e^{h_2},\qquad U(t)=\sum_{n\geq0}p_nt^n
       =\frac1{1-F(t)},\qquad F(t)=\sum_{n\geq1}f_nt^n.
\]
Identify the compact dual $\widehat{\mathcal O}=\operatorname{Hom}(\mathcal O,\R/\Z)$ with the $d$-torus using the integral power basis.
Its Haar probability $H$ is normalized Lebesgue measure in these coordinates; this is the Pontryagin-dual convention used in \citet{DGB16}.
On this group, write $(T\theta)(x)=\theta(\alpha x)$ and $w(\theta)=\cos^2(\pi\theta(1))$.
The products $W_n=\prod_{j=0}^{n-1}w\circ T^j$ and $V_n=\prod_{j=1}^{n}w\circ T^{-j}$ both integrate to $p_n$.
Under the additional hypothesis $U(R)<\infty$, define the probability densities
\[
 u_R=U(R)^{-1}\sum R^nW_n,\quad
 v_R=U(R)^{-1}\sum R^nV_n,\quad C_R=u_R*v_R.
\]
The convolution is already in $L^1(H)$ by transience.

\begin{theorem}\label{ret:critical-regularity}
If, in addition,
\begin{equation}\label{ret:fractional-premise}
 \sum_{n\geq1} n^aR^np_n<\infty\qquad(0<a<1),
\end{equation}
then $C_R$ belongs to weak $L^{1/(1-a)}$ and to every strong $L^p$ with $1\leq p<1/(1-a)$.
Weak $L^q$ means $\sup_{y>0}y^qH\{|C_R|>y\}<\infty$; the strong endpoint is not asserted.
\end{theorem}
\begin{proof}[Proof outline]
For $\delta>0$, the complete Abel deficit $U(R)-U(Re^{-\delta})$ is $O(\delta^a)$, giving an $O(\delta^a)$ error in $L^1$ and an $O(\delta^{a-1})$ bound for the one-sided subcritical convolution.
Balancing these estimates proves the distribution bound; dyadic differences give the $L^p$ approximation rate $O(\delta^{a-1+1/p})$.
\end{proof}

\begin{proposition}\label{ret:first-return-moment}
Writing $\eta_R=1/U(R)$, the complete first-return moment is equivalent:
\[
 \sum n^aR^nf_n\ \leq\ \sum n^aR^np_n
       \ \leq\ \eta_R^{-2}\sum n^aR^nf_n.
\]
\end{proposition}
\begin{proof}
The renewal identity gives $p_n\ge f_n$ and $U=\sum_{k\ge0}F^k$.
Since $0<a<1$, $(n_1+\cdots+n_k)^a\le\sum_jn_j^a$.
Applying this to each convolution coefficient bounds the complete moment by $\sum_{k\ge1}kF(R)^{k-1}\sum_n n^aR^nf_n$.
Here $1-F(R)=\eta_R>0$, so the first sum is $\eta_R^{-2}$.
\end{proof}

These are conditional statements about the entire return law.
Neither critical transience nor \eqref{ret:fractional-premise} is known here for an arbitrary infinite nontorsion core.
The torsion control $\alpha^4=-1$ illustrates why the distinction matters: its critical convolution is unbounded but satisfies $\int e^{C_R/64}\,dH\leq4e^{17/32}$ and hence every finite $L^p$ bound.
Infinite critical mean alone therefore does not imply poor finite-$p$ regularity, and finite-$p$ regularity alone does not resolve the gap.

\paragraph{The remaining Fourier input.}
A complete fractional return moment supplies the regularity just stated, but its existence on an arbitrary nontorsion infinite core remains an additional mathematical obligation.
A useful finite-window witness must retain its frequency separation, multiplier and smoothing costs.
In particular, varying the original input must not erase the uniform scale required by the gap question.

\subsection{Neutral completion, graph compression and finite capture}
\label{sec:structures}

This subsidiary strategy asks whether an actual integer or graph construction can retain the original entropy while removing the unbounded complexity that prevents finite capture.
A local neutral block is a possible building element; all connections to the rest of the construction must be paid.
The following complete families make that distinction concrete and also inform the determinant strategy in Section~\ref{sec:determinants}.

\subsubsection{A concrete presentation and its absolute companion}
\label{str:conventions}

\begin{definition}\label{str:proper-presentation}
Let $V$ be an integral upper triangular matrix with every diagonal entry equal to one; we call this a proper presentation.
Set
\[
 p_V(t)=\det(tV+V^T),\qquad
 V^-_{ii}=1,\quad V^-_{ij}=-|V_{ij}|\quad(i<j).
\]
The logarithmic measure $m(p_V)$ uses every factor of this monic reciprocal polynomial.
For later cost comparisons write $E(V)=\sum_{i<j}V_{ij}^2$ and $H(V)=\max_{i,j}|V_{ij}|$.
\end{definition}
The literal signed two-sheet construction $W(V)$ doubles the rank and is again integral upper triangular with diagonal one.
To specify it, interlace two copies of each ordered vertex.
Replace an entry $v=V_{ij}$ by the $2\times2$ block
\[
 \begin{pmatrix}(v-|v|)/2&-(v+|v|)/2\\
                 -(v+|v|)/2&(v-|v|)/2\end{pmatrix}.
\]
\begin{proposition}\label{str:signed-cover}
On the symmetric and antisymmetric subspaces this block acts by $-|v|$ and $v$, respectively.
Thus
\begin{equation}\label{str:cover}
 p_{W(V)}=p_Vp_{V^-},\qquad
 m(p_{W(V)})=m(p_V)+m(p_{V^-}),\qquad E(W(V))=2E(V).
\end{equation}
\end{proposition}
The proposition follows by decomposing the two sheets into their symmetric and antisymmetric parts.
The companion is a genuine extra expense.
Neutral factors have zero Mahler measure but remain part of the matrix identity.

\subsubsection{A family where every repeated local change has linear cost}
\label{str:repeated}

For $k\geq3$ use the ordered vertices
\[
 (a_0,b_0,c_0),\ldots,(a_{k-1},b_{k-1},c_{k-1}),
 z_1,\ldots,z_k,a_k.
\]
Define $V_k$ by diagonal entries one, the diamond entries
\[
 V_k(a_i,b_i)=V_k(a_i,c_i)=V_k(b_i,a_{i+1})=-1,
 \qquad V_k(c_i,a_{i+1})=1,
\]
the corner $V_k(a_0,a_k)=-1$, and entries $-1$ along $a_0,z_1,\ldots,z_k,a_k$.
All remaining strict entries are zero.
\begin{proposition}\label{str:base-family}
This is a connected height-one presentation of rank $4k+1$ and energy $5k+2$.
Its complete scalar is
\begin{equation}\label{str:original}
 p_{V_k}(t)=(t+1)^{k-1}(t-1)^{2k-2}F_k(t),\qquad
 F_k(t)=\frac{t^{k+1}A(t)-A^*(t)}{t-1},
\end{equation}
where $A=t^4-4t^3-2t^2+1$ and $A^*=t^4-2t^2-4t+1$.
The quotient is a polynomial identity, including $t=1$.
The unique exterior root $\lambda_k$ of $F_k$ lies in $(4,5)$, and
\begin{equation}\label{str:original-cost}
 m(p_{V_k})=\log\lambda_k\in(\log4,\log5).
\end{equation}
\end{proposition}
Thus the varying positive factor does not supply a sequence of measures tending to zero.

Choose any subset $S\subseteq\{0,\ldots,k-1\}$.
In each selected cell replace $(a_i,b_i,c_i)$ by $(b_i,c_i,a_i+b_i+c_i)$, using
\[
 Q_3=\begin{pmatrix}0&0&1\\1&0&1\\0&1&1\end{pmatrix}.
\]
The resulting block-diagonal determinant-one matrix $Q_S$ gives $X_{k,S}=Q_S^TV_kQ_S$.
Every intercell and endpoint attachment is transformed along with the cell.
The result is still an upper unitriangular integral matrix of the same rank, height and energy, with the entire scalar \eqref{str:original} unchanged.

\begin{theorem}\label{str:uniform-cost}
Uniformly over every subset $S$, including all cells and all growing cardinalities,
\[
 \frac{m(p_{X_{k,S}^-})}{k}\longrightarrow\log2,
 \qquad
 \frac{m(p_{W(X_{k,S})})}{k}\longrightarrow\log2.
\]
For every $k\geq3$ one also has $m(p_{X_{k,S}^-})\leq k\log2+\log(20k^2+4k)$.
\end{theorem}

The complete companion depends only on $|S|$.
Eliminating the two middle vertices of a cell leaves the same matching recurrence
\[
 T_{-1}=0,\quad T_0=1,\quad T_1=(t-1)^2,\qquad
 T_j=(t-1)^2T_{j-1}-4t^2T_{j-2}.
\]
Selection changes the two oriented cycle exponents, but not this bulk recurrence.
Proving the theorem for a growing number of selected cells requires the full circle integral, rather than finitely many leading coefficients.
On each circle $|t|=r>1$, the two recurrence roots have different moduli away from the real axis.
Comparing the dominant root with all cycle terms, including possible exponential ties, gives the lower exponential rate.
The negative logarithmic parts are uniformly integrable because, for any angular set of measure $\delta$, each linear factor contributes at most $\delta[1+\log(2/\delta)]$ to its negative integral.
The exact Mahler integral \cite{Smy08} then pays the whole radial expense $(4k+1)\log r$ before $r$ decreases to one.
This establishes all subset quantifiers; it does not establish optimality over other integer bases of the same congruence class.

\subsubsection{Changing the bulk: neutral cells still have costly couplings}
\label{str:changed}

A different actual basis does change the recurrence.
In every cell set
\[
 (u_i,r_i,w_i)=(c_i,a_i+b_i+c_i,-a_i-c_i),\qquad
 R_3=\begin{pmatrix}0&1&-1\\0&1&0\\1&1&-1\end{pmatrix},
 \quad\det R_3=1.
\]
Let $Y_k=Q^TV_kQ$ with these blocks and identity on the remaining vertices.
Its internal block is $\left(\begin{smallmatrix}1&1&-1\\0&1&-1\\0&0&1\end{smallmatrix}\right)$; its adjacent block is $(1,0,-1)^T(0,1,-1)$.
The transformed initial attachments have row $-(0,1,-1)$, the final attachment has column $(1,0,-1)^T$, and the $z$ path is unchanged.
These specify all couplings.
The matrix is proper, has height one, rank $4k+1$, and energy $8k+2$, paying an increase of $3k$.
Its original polynomial remains \eqref{str:original}.

Write $U_k=Y_k^-$.
Each isolated absolute cell is
\[
 T_3=\begin{pmatrix}1&-1&-1\\0&1&-1\\0&0&1\end{pmatrix},
 \qquad\det(tT_3+T_3^T)=(t+1)(t-1)^2.
\]
Although this cell is neutral, it does not split from its neighbors.
Put $u=(1,0,1)^T$, $v=(0,1,1)^T$ and $B=tT_3+T_3^T$.
The open $j$-cell pencil $K_j$ has diagonal blocks $B$, upper blocks $-tuv^T$ and lower blocks $-vu^T$.
Its determinant is $(t+1)^jh_j$, where
\begin{equation}\label{str:changed-recurrence}
 h_0=1,\quad h_1=(t-1)^2,\qquad
 h_j=(t^2-6t+1)h_{j-1}
       -4t(2t+1)(t+2)h_{j-2}.
\end{equation}
This genuinely differs from the previous recurrence.
The open chain, however, is not the complete companion.

For an explicit closed formula put $C_j=\sum_{h=0}^jt^h$, $b=tC_{k-1}/C_k$, and let $f$ be $v$ in the first cell and zero in the others, while $g$ is $u$ in the last cell.
Schur elimination of the whole endpoint path \cite[Appendix~A.5.5]{BV04} gives
\begin{equation}\label{str:closed}
 p_{U_k}=C_k\det
 \begin{pmatrix}
 K_k-bff^T&-tg-(t+t^{k+1}/C_k)f\\
 -g^T-(1+1/C_k)f^T&t+1-b
 \end{pmatrix}.
\end{equation}
This is a full polynomial identity after clearing denominators.
Its factor $C_k$, corner terms, and exceptional denominator roots remain; no claim that the open determinant divides $p_{U_k}$ is used.

\begin{proposition}\label{str:changed-cost}
For every $k\geq3$,
\[
 \Tr(-U_k^{-1}U_k^T)=2^k+13k,\qquad
 m(p_{U_k})\geq\log(2^k+9k+2)>k\log2.
\]
Moreover,
\[
 m(p_{U_k})\leq\tfrac12\bigl[(k-1)\log6+2\log7
                   +(2k-1)\log8+2\log9+(k-1)\log10\bigr].
\]
The full cover adds $\log\lambda_k$ to these companion costs and has rank $8k+2$ and energy $16k+4$.
\end{proposition}

The trace identity counts every alternative increasing path between the endpoints of each existing edge.
The long connections contribute the term $2^k$; omitting them would erase the obstruction.
For a real monic reciprocal polynomial of odd degree $n$, remove one forced root $-1$ for the estimate and pair all other roots reciprocally.
Their real trace is at most $n-4+2\cosh m$.
With $n=4k+1$, the exact trace therefore gives $e^m\geq2^k+9k+2$.
The upper bound follows from the complete support-degree histogram and Hadamard's inequality integrated on the unit circle.

The useful conclusion is specific: this changed basis makes every isolated absolute cell neutral, yet its actual couplings still cost linearly in $k$.
Neither this family nor the repeated-cell family optimizes over all proper integral bases, all completions, or all scalar targets.
A new structural construction would have to control the full closed determinant and its attachments, while preserving the scalar return to the unrestricted Mahler problem.

\paragraph{The remaining completion input.}
A successful structural construction must retain the original scalar factor, all complementary factors, the integral realization and the chosen trace normalization.
The repeated-cell results and the changed-basis example show that neutral local cells can retain a linear full coupling cost.
They do not optimize over every integral change of basis or every possible completion.


For finite-set attainment, Proposition~\ref{fo:primary} gives a useful test: lengthening one positive primary component does not create a finite-entropy nonattainment example.
A completion mechanism must handle the actual unrestricted flow or supply a new quantitative restriction on distinct positive weights.
